\section{讨论}
\label{sec:discussion}

\subsection{托管 Jev 是否适合作为 \pol{} 的治理工具？}
\label{sec:verdict}

前述各节以检测器与裁决者之分回答了研究问题。
此处我们逐项功能地说明这一回答对规范本身意味着什么，使该结论能够对照已编码的发现及其确定性（\cref{tab:keyfindings}）加以核查，而非靠推断得出。
\Cref{tab:verdict} 列出类型化模型可能被要求承担的每一项 \pol{} 功能，指明其所涉及的要求，给出决定该功能结论的一两项发现，并给出三种结论之一：\emph{证据不支持}，即已编码证据显示该功能失效，或该功能依赖于证据所否定的某一属性；\emph{仅作有条件的检测器}，即该功能可以在 \cref{sec:design} 的设计之内、以 \cref{sec:intro} 所述的检测器角色、作为留有记录的检测器履行，而证据不允许以其他任何方式履行；\emph{未经测试}，即没有已编码研究测量该条款在 \pol{} 构念上所要求的内容。
每一个确定性条目均为相应发现在 \cref{tab:keyfindings} 中的评级；没有一项为中等，凡依赖有争议发现的结论均作了相应标记。
这些结论针对的是 2026 年 9 月 15 日至 10 月 8 日各版本的托管 \jev{}，并在发现的适用范围为“两者”之处，针对一般类型化模型；它们不是对后续模型的预测。

{\scriptsize
\setlength{\tabcolsep}{3pt}
\begin{longtable}{@{}>{\raggedright\arraybackslash}p{0.12\linewidth}>{\raggedright\arraybackslash}p{0.04\linewidth}>{\raggedright\arraybackslash}p{0.34\linewidth}>{\raggedright\arraybackslash}p{0.14\linewidth}>{\raggedright\arraybackslash}p{0.10\linewidth}>{\raggedright\arraybackslash}p{0.17\linewidth}@{}}
\caption{逐项功能的 \pol{} 结论。\emph{要求}：所涉及的 \cref{tab:clauses} 中的要求。\emph{证据}：决定该结论的已编码发现及其决定性数字。\emph{结论}：NS，证据不支持；DC，仅作有条件的检测器，且须留有记录（即 \cref{sec:design} 的设计）；NT，未在 \pol{} 的构念上经过测试。同一单元格可对某项功能的一部分给出 NS、对另一部分给出 NT；“审慎原则”标记的是依据测量缺失、而非依据测得的失效所作的限制（\cref{sec:scope-certainty} 的规则）。\emph{确定性}：决定性发现在 \cref{tab:keyfindings} 中的评级；$^{c}$~有争议；“未评级”标记并非关键发现、且依赖单项研究的结果。所有评级均为低或极低。}
\label{tab:verdict}\\
\toprule
\pol{} 功能 & 要求 & 证据显示的内容 & 结论 & 确定性 & 尚需什么 \\
\midrule
\endfirsthead
\multicolumn{6}{@{}l}{\emph{\Cref{tab:verdict}（续）。}}\\
\toprule
\pol{} 功能 & 要求 & 证据显示的内容 & 结论 & 确定性 & 尚需什么 \\
\midrule
\endhead
\bottomrule
\endlastfoot
安全阀：阻断输入中的恨语（\clause{5.4.1}） & G1, G2, G3, G6 &
零样本下对有害英文内容的排序效果良好（中位 AUROC 0.886，每轮成本 \$0.30，对照为 \$18.96~\cite{arx29429}），在 HateCheck 诊断集上与一个前沿 LLM 的 macro-F1 相差不超过 1.6 个点，成本低 97\%~\cite{arx03324}；但固定阈值失效（阈值 0.5 时 F1 为 0.158，拟合阈值下为 0.947~\cite{jkf87_2026replication}），一种针对 \jev{} 的注入与一种通用后缀达到 97--100\% 的攻击成功率~\cite{arx04985}，一条插入的观点翻转了 12.1\% 的决策~\cite{arx31142}，互换选项定义改变了 32.5--50.5\% 的答案~\cite{arx26758,arx00346}。 &
DC：仅在独立构建的检测器取得一致时自动\emph{暂扣}；最终阻断由人作出 &
低（排序；操纵；选项名称）；低$^{c}$（阈值） &
一个中文三值基准；在本地留出数据上认证并随时间重新认证的阈值；每类攻击的红队成功率低于预先设定的目标；通过重命名测试的中性标识符 \\
\addlinespace
输出抑制：对 PAI 自身输出的恨语截断（\clause{4.3.3.1}） & G1, G3, G5 &
所有系统对仇恨语言与仅具冒犯性语言的区分都很弱（三分类 macro-F1：\jev{} 0.532，前沿 LLM 0.604，有监督 TF-IDF 0.639~\cite{arx03324}）；在没有检索先例的情况下，准确选出被违反的仇恨言论政策的比例为 29.4\%~\cite{arx07953}；错误集中于总体指标所掩盖的子类型（167 次漏检中有 130 次的置信度 $\geq 0.95$~\cite{arx33401}）。 &
DC：自动\emph{修复} AI 自身的输出是可逆的，且处于检测器角色之内；截断记录必须显示被删除的内容 &
低（子类型；选项名称）；低（排序） &
以中文测得的恨语相对于冒犯性语言及不在场状态的召回率；子类型召回率下限；每次截断的记录，向其所针对的人开放 \\
\addlinespace
实时矫正支持：识别与矫正建议（\clause{4.3.3.2}） & G1, G5 &
识别见上一行。类型化决策只返回一个概率，别无其他~\cite{willison2026jev}，因此矫正\emph{建议}需要一个生成式组件，而其依据并非该判定的依据（\cref{sec:g5}）；在中文情绪分类上，所有系统均表现较弱（macro-F1：\jev{} 21.06，生成式模型 19.82--23.31~\cite{arx08829}）。 &
NT：类型化接口仅涵盖识别；建议未经测试 &
低（排序）；未评级（中文情绪，单次运行） &
一项中文测试，检验在类型化判定之后生成的建议是否反映了产生该判定的因素 \\
\addlinespace
异常检测与风险预警（\art{7}(2)--(3)） & G1, G2, G6 &
在文本上，是附带上述条件的检测器；在非文本信号上，零样本 \jev{} 在网络流量特征上的得分低于多数类基线（9.8\% 对 22.8\%~\cite{arx00376}），在工业故障检测中将 57--94\% 的无故障窗口标记为异常，且在分布偏移下并不优于将每个窗口都标记为异常~\cite{arx09937}；伪装为正常运行的攻击召回率为 0.34，被一致性门控接受为良性的 76 个窗口中有 21 个是攻击~\cite{arx02048}。 &
文本上为 DC；在没有经训练模型的情况下，对数值或结构化信号为 NS &
低（子类型；操纵）；极低（一致性门控） &
以规范自身术语给出的“异常模式”定义；漂移与伪装攻击测试；以经训练的或基于规则的检测器作为独立的第二道核查 \\
\addlinespace
对齐的伦理验证（\art{6}(3)） & G3, G4, G5 &
在无人设时，\jev{} 回答价值问题的方式与其美国人设相似，沙特人设在英文中重现了人类沙特--美国差异的 87\%，在阿拉伯文中则为 62\%~\cite{arx36399}；在道德两难问题上，改写措辞对其答案的影响远大于重复运行~\cite{zen23023694}\zmark{}；在标注者分歧最大之处，其校准误差升至 0.339~\cite{zen23032384}\zmark{}，并且在人们以 55 比 45 分歧的条目上报告了 0.88~\cite{zen23179064}\zmark{}。 &
作为自动验证器为 NS，因为它依赖于中性的默认立场以及在人们存在分歧之处的校准，而两者均被证据否定；EAP 构念本身为 NT &
极低（文化）；低（校准是局部的） &
一个由价值负载的 \pol{} 问题构成、采用顾及分歧标签的基准，报告其无人设时的默认立场，并将有争议条目交由人处理 \\
\addlinespace
无需人工干预的争议裁决（\art{10}，\art{9}(7)） & G4, G5, G6 &
低成本同类评估器在 96.0\% 的判定中重复了 \jev{} 置信度最高的错误（上界）~\cite{arx29769}，并在 130 次高置信漏检中至多发现 5 次~\cite{arx33401}；明确的“不知道”选项在其为正确答案的条目中有 10.5\% 被选中，“以上皆非”（None）为 7--54\%~\cite{arx34024,arx39496}；通过两种接口提出的同一决策在 32.8\% 的情况下导致不同行动~\cite{arx37470}；预先固定的阈值在 14 项任务中的 8 项上未达到其准确率目标~\cite{arx24574}，一项预注册研究使这一发现有争议~\cite{arx33689}；输出不附带理由~\cite{willison2026jev}。 &
作为裁决者为 NS；仅对申请的排序与分诊为 DC &
低（同类模型错误；“不知道”）；低$^{c}$（预先设定的阈值） &
下文“何种情况会改变结论”中的四项条件，在一个中文 \pol{} 基准上，于一项经预注册并有独立重复的研究或两项独立研究中得到满足 \\
\addlinespace
对爱语与恨语的逐人量化（\clause{5.6}） & G3, G7 &
在一项非正式的强制选择测试中存在群体偏倚（“White”占首选的 51\%~\cite{simmons2026saidno}\gmark{}）；类型化与生成式模型在低资源语言上均出现性能退化~\cite{arx37647}；第三方插入的观点翻转了 12.1\% 的决策~\cite{arx31142}，因而汇总结果可能被操纵；没有研究测量行为判断中按说话者群体区分的差异化错误。 &
作为决策输入为 NS（操纵，低）；按群体的错误为 NT：在按群体与语言的错误公布之前，只对事件与行为评分，绝不对人评分（审慎原则，\cref{sec:scope-certainty}） &
极低（文化）；低（操纵）；未评级（群体偏倚，单项非正式测试；语言退化，单一基准） &
以中文公布的按群体与语言划分的错误率；同意、查阅本人记录与申诉；对照社会评分禁令的法律解读（\cref{sec:synthesis}，T5） \\
\addlinespace
公共决策中的评估维度（\clause{5.6}） & G7 &
总体规模的解读可以诚实地完成：一个人类--模型共享误差项在名义 80\% 与 90\% 水平下给出了 0.93 与 0.96 的回顾性覆盖率~\cite{arx27535}，经审计的概率样本给出了校正后的全州计数，但仅在按因素进行人工再校准之后~\cite{arx00213,arx10321}；模拟公众偏向某一人群~\cite{arx36399}，并低估分歧~\cite{zen23179064}\zmark{}。 &
DC：作为经审计的汇总测量，其测得误差须带入每一项结论；本身不作为决策输入 &
极低（文化）；未评级（审计设计，单项研究） &
一个由真实 \pol{} 条目构成、顾及分歧的审计样本；每一项估计都报告其与人之间的测得差距 \\
\end{longtable}
}

\paragraph{结论。}
按截至 2026 年 10 月 8 日的证据，托管 \jev{} 不适合承担 \pol{} 中任何“做决定”的功能。
这包括 \art{10} 下的争议裁决、\art{6}(3) 下的自动伦理验证、\clause{5.6} 的逐人量化（依据在于操纵，以及任何群体错误测量的缺失），以及安全阀的任何如下用法：对某人言论的阻断无须经人确认即成为最终结果。
决定性的发现是否定性的，并且在其低或极低的确定性范围内相互一致：在一种情境中拟合的阈值不能迁移到另一种情境（低，有争议）；决策可被模型所审计对象内部的内容操纵（低；定制注入的攻击成功率为 97--100\%~\cite{arx04985}，一条插入的观点翻转了 12.1\% 的决策~\cite{arx31142}）；选项名称可能压过其定义（低）；明确的“不知道”选项在其为正确答案时很少被选择，“以上皆非”的选择则不一致（低）；而此类模型会升级至的低成本评估器重复了其几乎全部高置信错误（低，上界）。
具有这些特性的裁决者，将依据可被对手操纵、取决于问题措辞、且很可能没有第二个自动意见能够发现的判定来限制言论、分配福利；\art{9}(7) 使人的监督不得干预，这将移除证据表明所必需的唯一一道核查。
所测试的类型化模型无一能逃脱这一结论，与之一同测试的低成本生成式评估器也未被证明能够逃脱：在进行了比较之处，类型化模型并非普遍更差（低，依据判断），但大多数失效从未进行比较，因此该结论并非专属于类型化模型，并且在生成式评估器于同一协议下接受检验之前，出于审慎原则适用于这一成本层级的自动裁决者（\cref{sec:scope-certainty}）。

对于检测而非决定的功能，即 \clause{5.4.1} 的安全阀、\clause{4.3.3.1} 的截断，以及 \art{7} 在文本上的异常检测（不包括在没有经训练模型时对数值或结构化运行信号的检测，\cref{tab:verdict}），托管 \jev{} 可以使用，但只能在 \cref{sec:design} 的设计之内使用；\clause{5.6} 的公共决策评估仅可作为经审计的汇总测量使用，并须将其测得误差带入每一项结论；支持这一点的正面证据是真实的，但附有条件。
它以生成式评估器成本的一小部分对有害英文内容作出良好排序（低），在 HateCheck 诊断集上与前沿模型的 macro-F1 相差不超过 1.6 个点，成本低 97\%~\cite{arx03324}；它在与其训练相似的任务上往往校准良好（低）；基于置信度的暂缓判定在分布内有效~\cite{arx00381,zen22885291}；在一项预注册研究中，仅在一条独立构建的规则同意时才接受其低风险判定，并未损失 macro-F1，尽管伪装攻击得以通过（极低）。
每一项正面证据都附有将其转化为设计要素的条件：排序在成为决定之前需要经本地拟合并重新认证的阈值；校准只在测得之处成立；暂缓判定只在分布内成立；一致性门控需要一个构造不同的第二检测器。
即便这一检测器结论也是一种外推：它依据的是二值仇恨言论的英文基准，而 \pol{} 要求以中文对爱语、恨语与不在场状态作出三值判断，且没有已编码研究测量这一构念；唯一的中文情绪结果显示所有系统均较弱（macro-F1 约为 21~\cite{arx08829}），仅有的中文决策基准涉及服务路由、诈骗检测与船舶设计工程，而非恨语~\cite{qin2026zhdecisionbench,arx09896}；而仇恨语言与仅具冒犯性语言的区分，即 \clause{1.5} 最需要的区分，在所有受测系统中都很弱~\cite{arx03324}。
因此，我们在陈述结论时附上其确定性：上述每一项发现均为低或极低，四项有争议，大多数依赖单一团队在模型面世最初几周内发布的预印本，而一项针对该规范自身构念、设计良好的研究即可能改变这些评级。

另有两项发现与 \pol{} 的关联超出任何单一功能。
类型化保障机制也可能充当可查询的判定接口（oracle）：在输入中植入条件指令后，仅凭托管 \jev{} 的标签，就在每一个测试案例中揭示了隐藏的性别、年龄组、心脏病与族裔属性，每个状态需两到六次查询（开源模型的标签未揭示任何属性）；此外，其判定提高了有害计划被评判的有害程度，其标签以 200 个标签训练出一个一致率达 92\% 的替身模型~\cite{arx04985}。这支持不向不受信任的调用方提供安全阀的概率，并将决策不需要的敏感字段排除在输入之外；对托管 \jev{} 而言，后者是这些措施中唯一能够消除该信号的一项，因为该攻击对获授权执行该任务的用户同样有效，且仅需标签；将不受信任调用方的查询限制在获授权的任务之内，可防止安全阀成为回答诸如有害计划可行性等问题的通用判定接口，而日志记录可使通过获授权任务进行的替身模型提取变得可见，但两者都不能阻止这种提取（\cref{sec:design}，要素~4）。这些措施均未经测试。
另外，一个封闭的托管模型无法按其字面要求满足 \art{4} 与 \art{5}(4)（T3）：其决策记录可以公开，其推理则不能；而与开源模型测得的鲁棒性之间的权衡（开源模型表现出最严重的名称反转~\cite{arx26758}）必须明确作出。

\paragraph{可借鉴之处。}
有若干经验无论使用何种模型都适用于 \pol{}，因为每一条都源自一项关乎接口或任务、而非关乎 \jev{} 的发现。
\begin{itemize}
  \item \emph{提出三值问题，并单独询问充分性。} 是/否问题会丢弃不在场状态（\clause{4.3.2}）；明确的“不知道”与“以上皆非”选项的选择并不一致（低），但在一项研究中，单独就证据是否充分提出的是/否问题，比置信度更能发现信息缺失~\cite{arx01006}（极低）。三值判定已纳入该工作组的试点方案~\cite{polgov}（\cref{sec:design}，要素~3；T2）。
  \item \emph{将中性标识符绑定到定义，并通过重命名加以测试。} 对所有受测类型化模型，选项名称都压过了定义（低），而数字或随机字符串标识符使托管 \jev{} 接近其噪声下限~\cite{arx00346}；\clause{1.5} 指出其概念并非道德词汇，道德词汇应置于定义之中，而不应置于选项名称之中（要素~3；核查清单第~4 项）。
  \item \emph{记录每一项决策。} 类型化判定不附带理由（\cref{sec:g5}），因此记录（问题、选项集及定义、输入哈希、概率、模型与校准版本、所援引的依据、复核结果）是依 \art{9} 唯一可被审计或质询的对象，也是 \art{5}(5) 可验证的决策链路的唯一基础（要素~7）。
  \item \emph{在本地校准，并随时间重新验证。} 在别处拟合的阈值未能迁移（低，有争议），校准是局部的（低）；静态基准上的仇恨言论排名不能预测在带时间戳或经新词替换的数据上的排名~\cite{dibonaventura2025hatevolution}，因此认证有其日期与语言（要素~4；核查清单第 2、3、11 与 16 项）。
  \item \emph{仅在构造不同的检测器取得一致时才触发自动动作。} 同类模型重复了高置信错误（低）；仅在一条独立构建的规则同意时才接受判定，并未损失 macro-F1，但伪装攻击仍然通过（极低），因此独立性必须被构建而不能被假定，且须测量其漏检率（要素~3；核查清单第~10 项）。
  \item \emph{凡可能限制人的链条，均以人作为终点，并审计高置信决策。} 最具破坏性的错误是高置信且共有的~\cite{arx29769,arx33401}，且没有研究在治理条目上将类型化模型与受训审核者进行比较；在此类研究出现之前，带有盲法随机审计的人工终点是证据未予削弱的唯一核查（要素 5 与 6；T4，选项~(b)）。
  \item \emph{不向不受信任的调用方提供概率，将不需要的敏感字段排除在输入之外，限制并记录不受信任的查询，并将相互独立的标记拆分到不同的调用中。} 概率反馈帮助了攻击者~\cite{arx30243}，仅凭标签即泄露了属性并训练出替身模型~\cite{arx04985}，一条植入的注释在一次调用中改变了多个答案~\cite{zen23038928}\zmark{}（要素 3 与 4）。
  \item \emph{在按群体与语言的错误公布之前，只对事件与行为评分，绝不对人评分。} 这既源自证据，同样也源自 \clause{4.3.1}（T5）。
  \item \emph{首先构建构念基准。} 上述每一项结论都是从英文二值任务外推而来；一个经独立标注、顾及分歧、带时间戳、以中文标注爱语、恨语与不在场状态并设有留出集的基准，是将其中任何结论从外推转为测量的前提（见下文开放问题）。
\end{itemize}

\subsection{范围、确定性、局限性与开放问题}
\label{sec:scope-certainty}

\paragraph{超越单一规范。}
任何在 AI 系统之前设置自动过滤器、将争议交给 AI，或为公共目的对人的行为进行评分的提案，都面临 G1--G7 这些问题；并且，在下文关于确定性的保留之下，\cref{sec:design} 的设计结论很可能也适用于它。
\cref{sec:paradigms} 中的比较使这一点具体化：宪法分类器、中国关于停止违法生成的义务以及 DSA 关于自动化手段的条款，都在朝推理期拦截层的方向推进，随后各自都需要回答 G1--G5，并在其将监督自动化之处回答 G6 的升级问题，而其自身文本对这些问题至多只给出了部分回答（《生成式人工智能服务管理暂行办法》要求记录对用户采取的措施并设立投诉渠道，第 14 条第 2 款与第 15 条~\cite{cac2023genai}）；G7 与 G6 的其余部分只在某一提案同时将争议交给 AI 或对人评分时才会出现。
本案例的增益在于表明：一份治理文本可以逐条解读为一组可检验的要求，而这样做会暴露出在仅把文本当作价值宣言来阅读时无从察觉的张力。
在本案例中，该框架自身的若干承诺，即不在场状态、关于爱与恨并非道德标签的告诫，以及对人与行为的区分，与证据所推荐的做法相一致；张力出现在其他条款与之相拉扯之处。

\paragraph{确定性。}
\Cref{tab:keyfindings} 对摘要与结论中使用的每一项发现进行了评级。
这些规则计入与某项发现相关的每一项已编码研究，将有共同作者的研究计为一项，只承认其他团队基于自身数据或重新实现所作的研究为独立重复，按系统组计数证据，并仅依据托管模型证据对关于托管 \jev{} 或一般类型化模型的发现进行评级（开源模型研究仅作佐证），且当一项设计强度至少相当的研究指向相反方向时降低评级（\cref{sec:appraisal}）；每项已编码研究在各项发现中的归属均已发布（\nolinkurl{data/certainty_evidence.csv}），评级由 \nolinkurl{scripts/rate_certainty.py} 据此计算得出。
这些规则在第 12--15 轮审查期间、即更新研究完成编码之后，根据审查者的意见作了修订，所有发现均据此重新评级；这偏离了曾冻结主窗口评级的更新方案，因此该表在当前评级旁列出了 10 月 4 日给定的评级。
没有发现达到中等确定性。若开源模型研究可以提高评级，“插入内容会操纵判定”将为中等；而在我们的规则下，它们只起佐证作用。
有四项发现有争议：注入很少选中攻击者的目标，评级为极低；升级至更强的推理模型可在不损失准确率的情况下节省成本，评级为低；以及两项阈值发现，即固定阈值不可迁移、预先设定的阈值在新数据上达不到其错误目标，评级为低。这两项均得到一项预注册研究的支持，在该研究中，预先固定的阈值在 14 项任务中的 8 项上未达到其准确率目标~\cite{arx24574}，并得到若干其他 \jev{} 研究的支持；两项均受到一项关于问题类型变化的预注册研究的反对，在该研究中，\jev{} 预先设定的门控在 7 个模板中的 6 个上保持在其目标之内~\cite{arx33689}，第二项还受到两项阈值达到目标的更新研究的反对~\cite{zen22935043,zen23075657}\zmark{}；由于该预注册相反研究与最强的支持研究同级，两项均由中等降为低。
所有评级均为低或极低，原因在于来源是模型发布后数日内发布的预印本，许多结果仅依赖一项研究，对一项重点研究的唯一一次再现重跑的是其作者的代码与数据，且没有研究测量该规范自身的构念。

\paragraph{局限性。}
语料库涵盖十七天内的投稿，由未经同行评审的预印本与灰色文献组成，其中许多出自快速工作的单一作者；在一个社区追踪项目评级的 27 项早期研究中，没有一项被评为低偏倚风险~\cite{rafe2026awesome}。
关于失效的研究在模型面世的最初几周内易于撰写，因而可能被过度代表；Zenodo 记录比 arXiv 论文更常给出否定结果，尽管剔除它们对总体比例的影响很小（见\cref{sec:across}）。
若干研究使用的是开源模型而非托管模型，有些研究未报告托管模型的版本；我们每次都予以说明，但计票结果将托管模型与开源模型合并统计。
各项要求源自 \pol{}，因此本分析回答的是它的问题，可能遗漏类型化模型在其未提出的问题上的长处。
筛选、提取、编码与核验均由 AI 智能体在作者指导下完成，且两名编码者是同一模型家族的 AI 智能体，因此一致性统计量（$\kappa = \nKappa{}$，95\% CI \nKappaLo{}--\nKappaHi{}，样本以 Q 编码为主）是该编码手册对人类编码者可重复程度的上界；对全部更新配对的全量编码得到 $\kappa = \nUKappa{}$（95\% CI \nUKappaLo{}--\nUKappaHi{}），其区间与前者重叠，是对 AI 编码者而言更好的估计，这意味着主窗口中只编码一次的编码有一部分在第二次审读时会得出不同结论。
主计数涵盖截至 10 月 4 日检索到的记录，且主检索的查全率并不完整：在日期位于窗口内的 103 项已编码研究中，有 18 项仅由 10 月 8 日的更新检索发现（\cref{sec:update}）；更新检索的结果单独报告、未予合并，其中包括 \nUSetReScreened{} 条无法确认此前曾被排除或遗漏的 Zenodo 记录；10 月 8 日之后公布的论文未予纳入。
没有研究在治理任务上将类型化模型与受过训练的人类审核者进行比较，因此我们无法判断其错误是否比它所协助的人更糟。

\paragraph{何种情况会改变结论。}
如果在一个留出的、经独立标注的中文三值 \pol{} 基准上，一个采用风险控制阈值的类型化模型
（i）将其误阻断率与漏检率保持在部署责任方预先设定的目标之内，
（ii）在名称与定义互换下改变的答案少于其重复噪声下限的两倍，
（iii）在单一观点插入和 64 次查询的自适应攻击下翻转的决策少于 5\%，并且
（iv）与一个独立的受训审核者小组之间的一致程度，至少不低于这些审核者彼此之间的一致程度，
且上述结果见于一项经预注册并有独立重复的研究，或见于两项独立研究，我们将修正“非裁决者”的结论。
如果在这样的基准上，其 AUROC 低于 0.75，或两个独立构建的检测器相互重复了对方 80\% 以上的高置信错误，或在经认证的阈值下复核负荷超过了为之预算的审核者容量，我们将修正“检测器”的论断。
放宽一项保障机制所需的证据，应多于增设一项保障机制所需的证据，因为错误地移除一项保障机制会限制人的言论与资源，而错误地增设一项保障机制只耗费复核时间。
如果在同一基准上，受训审核者与独立小组的一致程度低于经门控的检测器集成，或在超过预设比例的审计案例中将其判定向模型所显示的判定靠拢，我们将修正“具有实质后果的决定最终由人作出”这一建议。
而如果在同一协议下检验的生成式评估器在类型化模型未通过之处通过检验，结论将变为专属于类型化模型；在此之前并非如此。

\paragraph{开放问题。}
\begin{itemize}
  \item \emph{一个 \pol{} 基准。} 目前没有公开数据集对爱语、恨语与不在场状态进行标注。一个经独立标注的中文基准，采用顾及分歧的标签~\cite{davani2022dealing}，参照现有中文冒犯性语言资源的模式构建~\cite{deng2022cold,lu2023toxicn}，并将留出集保存在任何代码仓库之外，是上述每一项建议的首要前提；该工作组公开的 20 个试点条目明确不作为排名集~\cite{polgov}。由于静态的仇恨言论基准会随语言变化而失去预测价值~\cite{dibonaventura2025hatevolution}，此类基准应标注时间戳并定期更新，包括含有新出现词语的配对条目。
  \item \emph{人类审核者与类型化检测器的比较}，在相同的治理条目上进行，包括审核者彼此之间的一致性及其受模型判定影响的程度。
  \item \emph{失效的对照比较。} 大多数否定性发现没有生成式对照；在加入对照的情况下重新运行这些研究，将显示哪些失效是类型化模型所特有的。
  \item \emph{可查询保障机制的鲁棒性。} 隐藏或粗化概率、添加噪声、限制调用频率、将查询限制在获授权的任务之内，或将敏感字段排除在输入之外，能否降低攻击成功率，以及会给合法用户带来多大代价，尚未得到测量。
  \item \emph{忠实的解释。} 由另一模型事后撰写的解释能否反映产生类型化判定的实际因素，需要专门评估，方能满足 \art{9}(2) 的要求。
  \item \emph{逐人度量。} 若实施 \clause{5.6}，则在计算或存储任何逐人度量之前，必须先测量跨群体与跨语言的差异化错误。
  \item \emph{重复研究。} 大多数发现依赖于单一团队的单一研究；基于新数据的独立重复研究（对 \emph{Just Ask Jev} 的再现重跑的是作者的代码与数据~\cite{jkf87_2026replication}）几乎对每一项要求而言都是最有价值的下一步贡献。
\end{itemize}
